\documentclass[11pt]{article}
\usepackage[margin=1in]{geometry}
\usepackage{amsmath,amssymb,amsthm,mathtools}
\usepackage[T1]{fontenc}
\usepackage[utf8]{inputenc}
\usepackage{enumitem}
\usepackage{booktabs,longtable,array}
\usepackage{hyperref}
\usepackage{xcolor}
\usepackage{microtype}
\hypersetup{colorlinks=true, linkcolor=blue!50!black, citecolor=blue!50!black, urlcolor=blue!50!black}

\newtheorem{theorem}{Theorem}[section]
\newtheorem{proposition}[theorem]{Proposition}
\newtheorem{lemma}[theorem]{Lemma}
\newtheorem{corollary}[theorem]{Corollary}
\newtheorem{claim}[theorem]{Claim}
\theoremstyle{definition}
\newtheorem{definition}[theorem]{Definition}
\newtheorem{assumption}[theorem]{Assumption}
\theoremstyle{remark}
\newtheorem{remark}[theorem]{Remark}

\newcommand{\KL}{\mathrm{KL}}
\newcommand{\Ent}{\mathrm{Ent}}
\newcommand{\Prob}{\mathsf{P}}
\newcommand{\E}{\mathbb{E}}
\newcommand{\R}{\mathbb{R}}
\newcommand{\N}{\mathbb{N}}
\newcommand{\Pcal}{\mathcal{P}}
\newcommand{\Mcal}{\mathcal{M}}
\newcommand{\Qcal}{\mathcal{Q}}
\newcommand{\Zcal}{\mathcal{Z}}
\newcommand{\Fcal}{\mathcal{F}}
\newcommand{\Gcal}{\mathcal{G}}
\newcommand{\W}{\mathcal{W}}
\newcommand{\one}{\mathbf{1}}
\newcommand{\dd}{\,d}
\newcommand{\supp}{\operatorname{supp}}
\newcommand{\argmin}{\operatorname*{arg\,min}}
\newcommand{\argmax}{\operatorname*{arg\,max}}
\newcommand{\esssup}{\operatorname*{ess\,sup}}
\newcommand{\ri}{\operatorname{ri}}
\newcommand{\dom}{\operatorname{dom}}
\newcommand{\Law}{\operatorname{Law}}
\newcommand{\Var}{\operatorname{Var}}
\newcommand{\Cov}{\operatorname{Cov}}
\newcommand{\diag}{\operatorname{diag}}
\newcommand{\Id}{\mathrm{Id}}
\newcommand{\Lip}{\operatorname{Lip}}
\newcommand{\osc}{\operatorname{osc}}
\newcommand{\diam}{\operatorname{diam}}
\newcommand{\proj}{\operatorname{proj}}
\newcommand{\e}{\mathrm{e}}

\newenvironment{argument}{\paragraph{Argument.}\begin{enumerate}[leftmargin=2em]}{\end{enumerate}}
\newenvironment{proofoutline}{\paragraph{Proof architecture.}\begin{enumerate}[leftmargin=2em]}{\end{enumerate}}


% Source-faithfulness apparatus used by the paper reconstructions.
\newcommand{\sourceanchor}[1]{\par\smallskip\noindent\textbf{Source anchor.} #1\par\smallskip}
\newcommand{\sourcecomment}[1]{\begin{remark}#1\end{remark}}
\newenvironment{sourceguide}
  {\begin{description}[style=nextline,leftmargin=3.2cm,labelwidth=3cm]}
  {\end{description}}

% Backward-compatible environments used by some of the older reconstructions.
\newenvironment{distilled}[1][]{\begin{theorem}[#1]}{\end{theorem}}
\newenvironment{distillation}{\begin{enumerate}[leftmargin=2em]}{\end{enumerate}}

\newenvironment{commentarynotes}{\begin{remark}\begin{enumerate}[leftmargin=2em]}{\end{enumerate}\end{remark}}

\title{Modernized Proof Notes for Fortet (1940): Resolution of Schrodinger's System}
\author{}
\date{Source PDF: \texttt{Fortet 1940 - JMPA\_1940\_9\_19\_1-4\_83\_0.pdf}}
\begin{document}
\maketitle

\section*{Source map and scope}
\begin{description}[style=nextline,leftmargin=3.2cm,labelwidth=3cm]
\item[Primary source] \texttt{\detokenize{Fortet 1940 - JMPA_1940_9_19_1-4_83_0.pdf}}
\item[Coverage] Fortet's existence and uniqueness proof for the Schrödinger system, including truncation, monotone iteration, mass identities, limiting arguments, and Bernstein's conclusion.
\item[Source anchors]
\begin{itemize}[leftmargin=2em]
  \item Theorems I--III
  \item successive approximation construction
  \item lemmas controlling truncation and limits
  \item final Bernstein application
\end{itemize}
\end{description}

The exposition below preserves the source proof order and proof technique.  Notation is normalized
and intermediate calculations are made explicit.  When the source contains a gap, ambiguity, or
apparently incorrect step, the reconstruction records the source step before adding a separate
commentary remark.

\section*{Reference and purpose}
This note expands Robert Fortet, ``Resolution d'un systeme d'equations de M. Schrodinger,'' \emph{Journal de mathematiques pures et appliquees}, 9e serie, tome 19, pp. 83--105, 1940 \cite{Fortet1940}.
The original is written in older French analytic language.  The reconstruction keeps the proof architecture, translates the symbols into modern measure-theoretic notation, and records the almost-everywhere conventions explicitly.

\paragraph{Modernization convention.}
Fortet writes two real intervals as $J'$ and $J''$ and writes the marginal densities as $\omega_1(x)$ and $\omega_2(y)$.  We keep this notation.  Integrals are Lebesgue integrals; ``measurable (B)'' means Borel measurable.  Equalities between functions are to be read outside null sets when the original proof says ``presque partout.''

\section{The Schrodinger system}
Fortet studies a nonnegative kernel $g(x,y)$ on $J'\times J''$ and two nonnegative marginal densities $\omega_1,\omega_2$.  His system is
\begin{align}
  \varphi(x) \int_{J''} g(x,y)\psi(y)\,dy &= \omega_1(x), \tag{S1}\label{eq:fortet-S1}\\
  \psi(y) \int_{J'} g(x,y)\varphi(x)\,dx &= \omega_2(y). \tag{S2}\label{eq:fortet-S2}
\end{align}
A solution is considered only up to the gauge
\[
  (\varphi,\psi) \sim (c\varphi,c^{-1}\psi), \qquad c\ne 0.
\]
Fortet calls a solution positive when $\varphi\ge0$ and $\psi\ge0$; he is mainly interested in such solutions.

\subsection*{Hypotheses}
The paper first lists basic hypotheses, then strengthens them to the working hypotheses used in the main existence and uniqueness arguments.
In modern terms, the working assumptions are:
\begin{enumerate}[leftmargin=2em]
\item $g$ is continuous and bounded above by a positive finite constant.
\item For each fixed $x$ or fixed $y$, $g(x,y)$ is not identically zero except possibly on a null exceptional set.
\item $\omega_1$ and $\omega_2$ are continuous nonnegative functions.
\item The total masses are normalized:
\[
  \int_{J'} \omega_1(x)\,dx = \int_{J''}\omega_2(y)\,dy = 1. \tag{N}\label{eq:fortet-normalization}
\]
\end{enumerate}
The probability interpretation of Schrodinger also often assumes $\int g(x,y)\,dy=1$, but Fortet's analytic argument for the system is stated for a broader normalized-marginal setting.

\section{Reduction to one scalar fixed-point equation}
From \eqref{eq:fortet-S2}, whenever the denominator is meaningful,
\[
  \psi(y)=\frac{\omega_2(y)}{\int_{J'}g(z,y)\varphi(z)\,dz}.
\]
Substituting this into \eqref{eq:fortet-S1} gives a single nonlinear equation for $\varphi$.
Fortet then changes variables by setting
\[
  h(x)=\frac{\omega_1(x)}{\varphi(x)}.
\]
Then $\varphi(x)=\omega_1(x)/h(x)$ and the single equation becomes
\begin{equation}
  h(x)=Q(h)(x)
  :=\int_{J''} g(x,y)
    \frac{\omega_2(y)}{G(h,y)}\,dy, \tag{F}\label{eq:fortet-fixed-point}
\end{equation}
where
\begin{equation}
  G(h,y):=\int_{J'}g(z,y)\frac{\omega_1(z)}{h(z)}\,dz. \tag{G}\label{eq:fortet-G}
\end{equation}
A nonzero, noninfinite solution of \eqref{eq:fortet-fixed-point} gives a solution of the original system by
\begin{equation}
  \varphi(x)=\frac{\omega_1(x)}{h(x)},
  \qquad
  \psi(y)=\frac{\omega_2(y)}{G(h,y)}. \tag{Recover}\label{eq:fortet-recover}
\end{equation}
This is the exact continuum ancestor of the alternating quotient updates in Sinkhorn scaling: divide one marginal by the current kernel-propagated scaling, then divide the other marginal by the propagated result.

\section{Fortet's class $(C)$}
Fortet introduces a class $(C)$ of functions on $J'$.
In modern notation, $H\in(C)$ if:
\begin{enumerate}[leftmargin=2em]
\item $H$ is Borel measurable;
\item $H$ is bounded below by a strictly positive constant;
\item $H$ is finite almost everywhere.
\end{enumerate}
The lower bound prevents the denominator in \eqref{eq:fortet-G} from blowing up by division by $H$.
The finiteness almost everywhere is enough because all equations are ultimately read modulo null sets on the supports of the marginals.

\section{The fundamental lemma}
Let
\[
  A_1:=\{x\in J': \omega_1(x)>0\}.
\]
For $H\in(C)$ Fortet forms $H'=Q(H)$ by \eqref{eq:fortet-fixed-point}.  The central lemma is the engine of the paper.

\begin{lemma}[Fortet's fixed-point comparison lemma, modern form]
If $H\in(C)$ and $H'=Q(H)$, then $H'$ is Borel measurable, is bounded below by a positive number on every finite subinterval of $J'$, and is finite almost everywhere on $A_1$.  Moreover, if either
\[
  H'\le H \quad\hbox{a.e.}
  \qquad\hbox{or}\qquad
  H'\ge H \quad\hbox{a.e.},
\]
then
\[
  H'=H \quad\hbox{a.e. on } A_1.
\]
\end{lemma}

\subsection{Proof of the regularity part}
Fortet uses finite truncations.  Choose increasing finite sets or finite intervals
\[
  J'_1\subset J'_2\subset\cdots\uparrow J',
  \qquad
  J''_1\subset J''_2\subset\cdots\uparrow J''.
\]
For $H\in(C)$ define truncated denominators
\[
  G_n(H,y)=\int_{J'_n} g(z,y)\frac{\omega_1(z)}{H(z)}\,dz.
\]
Because $g$ and $\omega_1$ are measurable and $H$ is positive and finite a.e., the integrand is nonnegative and measurable.  The functions $G_n(H,y)$ are well-defined and increase in $n$.
Let
\[
  G(H,y)=\lim_{n\to\infty}G_n(H,y).
\]
Fortet records that $G(H,\cdot)$ is of Baire class one, hence Borel measurable, positive on finite pieces in the required sense, and finite where the hypotheses force it to be finite.  The finite-piece lower bound comes from the nontriviality of $g$ on each finite interval together with the positive lower bound on $H$.

Then define truncated images
\[
  H'_n(x)=\int_{J''_n} g(x,y)\frac{\omega_2(y)}{G(H,y)}\,dy.
\]
The monotone convergence theorem gives
\[
  H'(x)=\lim_{n\to\infty}H'_n(x).
\]
Thus $H'$ is measurable.  Positivity and local lower bounds are inherited from the positivity/nontriviality of the kernel and marginal on finite intervals.  Finiteness a.e. on $A_1$ is not asserted directly from pointwise estimates; it is extracted from the mass identity below.

\subsection{The normalization identity}
Fortet defines the auxiliary nonnegative function
\[
  F(x,y)=g(x,y)\frac{\omega_2(y)}{G(H,y)}\frac{\omega_1(x)}{H(x)}.
\]
Truncating in both variables and applying Beppo-Levi/monotone convergence gives
\begin{align*}
  \int_{J'} \frac{H'(x)}{H(x)}\omega_1(x)\,dx
  &= \int_{J'}\left(\int_{J''} g(x,y)\frac{\omega_2(y)}{G(H,y)}\,dy\right)
       \frac{\omega_1(x)}{H(x)}\,dx \\
  &= \int_{J''}\frac{\omega_2(y)}{G(H,y)}
       \left(\int_{J'}g(x,y)\frac{\omega_1(x)}{H(x)}\,dx\right)dy \\
  &= \int_{J''}\omega_2(y)\,dy \\
  &=1. \tag{M}\label{eq:fortet-mass-identity}
\end{align*}
This identity is Fortet's equation (12) in modern notation.  It implies $H'$ is finite a.e. on $A_1$, because if $H'=\infty$ on a set of positive $\omega_1$-mass, the integral in \eqref{eq:fortet-mass-identity} would be infinite.

\subsection{The comparison conclusion}
If $H'\le H$ a.e., then
\[
  0\le \frac{H'}{H}\le 1 \quad\hbox{a.e.}
\]
and \eqref{eq:fortet-mass-identity} plus $\int\omega_1=1$ forces $H'/H=1$ $\omega_1$-a.e.  Therefore $H'=H$ on $A_1$ up to a null set.
Similarly, if $H'\ge H$ a.e., then $H'/H\ge1$ and the same integral identity forces equality.
This is the central fixed-point forcing step: one-sided comparison plus the mass identity collapses to equality.

\section{Successive approximation: the first existence theorem}
Fortet now constructs a sequence by alternating $Q$ with pointwise max/min truncations.  In modern notation, write pointwise maximum as $M(f,g)$ and pointwise minimum as $m(f,g)$.  The construction begins
\[
  H_1=1,
  \qquad H'_1=Q(H_1),
  \qquad H_2=m(H_1,H'_1),
\]
then
\[
  H'_2=Q(H_2),
  \qquad H'_3=Q(H'_2),
  \qquad H_3=M(H'_3,1/2),
\]
and continues with a clipped monotone scheme, arranged so that each $H_n$ remains in class $(C)$ and
\[
  H_{n+1}\le H_n,
  \qquad
  H'_n\le H'_{n+1}
\]
in the appropriate alternating sense.  The exact min/max choices are less important than the two invariants Fortet maintains:
\begin{enumerate}[leftmargin=2em]
\item every iterate used as input to $Q$ is in class $(C)$;
\item the output and input are eventually ordered, or else they generate monotone limits.
\end{enumerate}

\subsection{Case 1: an ordered iterate occurs}
Suppose that for some $n_0$,
\[
  H'_{n_0}\le H_{n_0}
  \quad\hbox{a.e.}
\]
Fortet cannot immediately apply the lemma if the relevant lower bound degenerates.  He therefore regularizes by
\[
  K_p=M(H'_{n_0},1/p).
\]
Then $K_p\in(C)$ and $K_p\downarrow H'_{n_0}$.  Set $K'_p=Q(K_p)$.  Monotone convergence shows that $K'_p$ decreases to a limit $K'$.  The comparison inequalities give
\[
  K'_p\le H_{n_0+1},
  \qquad
  K'\le H'_{n_0}.
\]
The mass identity applied to $K_p$ gives
\[
  \int_{J'}\frac{K'_p}{K_p}\omega_1\,dx=1.
\]
Passing to the limit gives
\[
  \int_{J'}\frac{K'}{H'_{n_0}}\omega_1\,dx=1.
\]
Since $K'\le H'_{n_0}$, the comparison part of the lemma forces
\[
  K'=H'_{n_0}\quad\hbox{a.e. on }A_1.
\]
Finally, the denominator identity shows $Q(K')=K'$, hence $K'$ is a positive solution of the fixed-point equation \eqref{eq:fortet-fixed-point}.  Formula \eqref{eq:fortet-recover} gives a positive solution of the Schrodinger system.

\subsection{Case 2: no ordered iterate occurs}
In the second case, no $H'_n$ falls below the corresponding $H_n$ on a set large enough to trigger Case 1.  Fortet takes monotone limits
\[
  H=\lim_n H_n,
  \qquad
  H'=\lim_n H'_n,
\]
and obtains
\[
  H\le H'. \tag{L}\label{eq:fortet-limit-order}
\]
He first proves that neither limit can vanish at an isolated point without vanishing a.e.; the positivity/nontriviality of $g$ and the recursive formula propagate vanishing through the integral equation.  The mass identity is then passed to the limit to obtain
\[
  \int_{J'} \frac{H'}{H}\omega_1\,dx=1.
\]
Together with \eqref{eq:fortet-limit-order}, this forces $H=H'$ a.e. on $A_1$, again by the comparison part of the lemma.  Thus the limit is a fixed point of $Q$, and \eqref{eq:fortet-recover} yields a positive solution of the system.

\subsection{Theorem I}
Fortet's first existence theorem can be modernized as follows.

\begin{distilled}[Fortet Theorem I, schematic modern form]
Assume the working hypotheses and the additional integrability/regularity hypotheses needed for the successive-approximation construction on the full intervals.  Then the Schrodinger system \eqref{eq:fortet-S1}--\eqref{eq:fortet-S2} has a positive solution $(\varphi,\psi)$, with $\varphi$ and $\psi$ recovered from a fixed point $h=Q(h)$ by \eqref{eq:fortet-recover}.
\end{distilled}

The proof is not a contraction proof.  It is an order/compactness proof: build monotone envelopes, use Beppo-Levi to pass to limits, and use the normalization identity to turn one-sided inequalities into equality.

\section{The class $(B)$ and the second existence theorem}
Fortet then isolates a condition on the kernel that implies the hypotheses needed in the first theorem locally and gives continuity of the recovered solution.
A function $g(x,y)$ is of class $(B)$ in $x$ on a finite interval $\mathcal F$ if convergence of integrals
\[
  \int g(x,y)f(y)\,dy
\]
which is finite a.e. on an open set containing $\mathcal F$ is in fact uniform for $x\in\mathcal F$.  The point of this definition is to upgrade almost-everywhere information to continuity in $x$.

Fortet gives sufficient conditions:
\begin{enumerate}[leftmargin=2em]
\item if $g$ is bounded below by a positive number independent of $x$ and $y$ on the finite window, then $g$ is class $(B)$;
\item more generally, if there is a uniform ratio bound
\[
  \frac{g(x_1,y)}{g(x_2,y)}<\beta
\]
for $x_1,x_2$ in the finite window and all relevant $y$, then $g$ is class $(B)$;
\item if $J'=J''=\R$ and $g(x,y)=U(x-y)$ with $U$ eventually monotone on both tails, then $g$ is class $(B)$ in both variables.
\end{enumerate}
The Gaussian kernel in Bernstein's case satisfies this tail monotonicity condition.

\subsection{Theorem II}
Let $A_1=\{x:\omega_1(x)>0\}$.  Fortet fixes a compact interval $\mathcal F\subset A_1$ and chooses an auxiliary positive continuous density $p(x)$ on $\mathcal F$ with
\[
  \int_{\mathcal F}p(x)\,dx=1,
\]
plus integrability inequalities making Theorem I applicable to the restricted system.  Solving the restricted system gives a local positive pair.  He then sets up a second fixed-point equation on $A_1$ and repeats the monotone argument.  The class $(B)$ property makes the kernel integrals uniformly convergent on compact subintervals of $A_1$, and therefore the recovered $h$ and hence $\varphi$ are continuous on $A_1$.

\begin{distilled}[Fortet Theorem II, modern form]
Under the working hypotheses, if $g$ is of class $(B)$ in $x$, then the system has a solution with
\begin{enumerate}[leftmargin=2em]
\item $\varphi(x)=0$ when $\omega_1(x)=0$;
\item $\varphi(x)>0$ and continuous when $\omega_1(x)>0$;
\item $\psi(y)\ge0$ is Borel measurable and vanishes, modulo null sets, only where $\omega_2(y)=0$.
\end{enumerate}
\end{distilled}

Fortet then notes that if $\omega_1>0$ on $A_1$ and $\omega_2>0$ on $A_2$, solving the restricted system on $A_1\times A_2$ and extending by zero outside those sets gives a solution on the original domains.

\section{Uniqueness}
Fortet proves uniqueness among positive Borel solutions, modulo the scalar gauge.
Let $(\varphi_1,\psi_1)$ and $(\varphi_2,\psi_2)$ be two positive solutions.  On $A_1=\{\omega_1>0\}$ set
\[
  h_i(x)=\frac{\omega_1(x)}{\varphi_i(x)}.
\]
Each $h_i$ satisfies the same fixed-point equation
\[
  h_i(x)=\int_{J''}g(x,y)
  \frac{\omega_2(y)}{G(h_i,y)}\,dy.
\]
Let
\[
  h(x)=M(h_1(x),h_2(x))=\max\{h_1(x),h_2(x)\}.
\]
Then $G(h,y)$ is finite a.e. on $A_2=\{\omega_2>0\}$, because $h\ge h_i$ implies $1/h\le1/h_i$ and the original denominators are finite.  Define
\[
  h'(x)=Q(h)(x)=\int g(x,y)\frac{\omega_2(y)}{G(h,y)}\,dy.
\]
The mass identity gives
\[
  \int\frac{h'}{h}\omega_1\,dx=1. \tag{U1}\label{eq:fortet-uniq-mass}
\]
Since $h\ge h_i$, we have $G(h,y)\le G(h_i,y)$ and therefore
\[
  Q(h)(x)=h'(x)\ge Q(h_i)(x)=h_i(x)
\]
for $i=1,2$.  Thus
\[
  h'\ge h.
\]
Together with \eqref{eq:fortet-uniq-mass}, the comparison lemma forces
\[
  h'=h \quad\hbox{a.e. on }A_1.
\]
From $h=h_1$ on the subset where $h_1\ge h_2$ and the equality of denominators, Fortet derives
\[
  G(h,y)=G(h_1,y)\quad\hbox{a.e. on }A_2,
\]
and hence the integrands force $h=h_1$ a.e.  Similarly $h=h_2$ a.e.  Therefore
\[
  h_1=h_2\quad\hbox{a.e. on }A_1.
\]
Undoing the substitution $h_i=\omega_1/\varphi_i$ gives $\varphi_1=\varphi_2$ after the chosen gauge normalization.  Without a fixed gauge, the result is uniqueness up to multiplication of one factor by a positive constant and the other by its reciprocal.

\begin{distilled}[Fortet Theorem III]
Under the working hypotheses, the Schrodinger system has at most one positive Borel solution modulo the scalar gauge $(\varphi,
\psi)\sim(c\varphi,c^{-1}\psi)$.
\end{distilled}

\section{Bernstein's Gaussian case}
Fortet's motivating special case, attributed to Bernstein, has
\[
  J'=J''=\R,
  \qquad
  g(x,y)=\frac{1}{a\sqrt\pi}\exp\left(-\frac{(x-y)^2}{a^2}\right),
  \qquad a>0,
\]
with continuous marginals.  The Gaussian kernel is a function of $x-y$ and is monotone on each tail, so it is class $(B)$ in both variables.  Therefore Theorems II and III imply:

\begin{distilled}[Fortet's Bernstein conclusion]
In the Gaussian Bernstein case, the Schrodinger system admits a positive continuous solution, and this solution is the unique positive Borel solution modulo the scalar gauge.
\end{distilled}

\section{Structural commentary}
\begin{commentarynotes}
\item The fixed-point operator  the continuum analogue of composing the two quotient updates in matrix scaling.
\item The mass identity \eqref{eq:fortet-mass-identity} is the device that removes residual scalar freedom from one-sided comparison arguments.
\item Fortet's proof is an order-and-compactness argument, not a Hilbert-metric contraction proof; its distinctive ingredients are truncation, monotone iteration, the mass identity, and local uniform compactness.
\item The finite-dimensional specialization retains the same algebraic skeleton: denominator map  update map  mass identity, and max--min comparison.
\end{commentarynotes}

\section*{Checklist of details retained}
\begin{itemize}[leftmargin=2em]
\item The two-equation system and gauge equivalence.
\item The reduction from $(\varphi,\psi)$ to a single equation $h=Q(h)$.
\item Class $(C)$ and why it is used.
\item The truncated-denominator construction and monotone convergence.
\item The normalization identity $\int Q(H)/H\,d\omega_1=1$.
\item The comparison lemma turning $Q(H)\le H$ or $Q(H)\ge H$ into equality.
\item The two-case successive-approximation proof.
\item The class $(B)$ local-uniform-convergence condition and the Gaussian/Bernstein application.
\item The uniqueness proof using $h=\max(h_1,h_2)$.
\end{itemize}

\begin{thebibliography}{99}
\bibitem{ChenGeorgiouPavon2021} Yongxin Chen, Tryphon T. Georgiou, and Michele Pavon. \emph{Stochastic Control Liaisons: Richard Sinkhorn Meets Gaspard Monge on a Schr\"odinger Bridge}. SIAM Review, 63(2):249--313, 2021. doi:10.1137/20M1339982. arXiv:2005.10963.
\bibitem{Leonard2014Survey} Christian L\'eonard. \emph{A Survey of the Schr\"odinger Problem and Some of Its Connections with Optimal Transport}. Discrete and Continuous Dynamical Systems--A, 34(4):1533--1574, 2014. doi:10.3934/dcds.2014.34.1533. arXiv:1308.0215.
\bibitem{Fortet1940} Robert Fortet. \emph{Resolution d'un systeme d'equations de M. Schrodinger}. Journal de mathematiques pures et appliquees, 9e serie, 19(1--4):83--105, 1940.
\bibitem{SinkhornKnopp1967} Richard Sinkhorn and Paul Knopp. \emph{Concerning Nonnegative Matrices and Doubly Stochastic Matrices}. Pacific Journal of Mathematics, 21(2):343--348, 1967.
\bibitem{FranklinLorenz1989} Joel Franklin and Jens Lorenz. \emph{On the Scaling of Multidimensional Matrices}. Linear Algebra and its Applications, 114/115:717--735, 1989. doi:10.1016/0024-3795(89)90490-4.
\bibitem{Carroll2004} Joseph E. Carroll. \emph{Birkhoff's Contraction Coefficient}. Linear Algebra and its Applications, 389:227--234, 2004. doi:10.1016/j.laa.2004.02.039.
\bibitem{EvesonNussbaum1995} Simon P. Eveson and Roger D. Nussbaum. \emph{An Elementary Proof of the Birkhoff--Hopf Theorem}. Mathematical Proceedings of the Cambridge Philosophical Society, 117(1):31--55, 1995.
\bibitem{BlanchetMurthy2019} Jose Blanchet and Karthyek R. A. Murthy. \emph{Quantifying Distributional Model Risk via Optimal Transport}. Mathematics of Operations Research, 44(2):565--600, 2019. doi:10.1287/moor.2018.0936. arXiv:1604.01446.
\bibitem{GaoKleywegt2023} Rui Gao and Anton J. Kleywegt. \emph{Distributionally Robust Stochastic Optimization with Wasserstein Distance}. Mathematics of Operations Research, 48(2):603--655, 2023. doi:10.1287/moor.2022.1275. arXiv:1604.02199.
\bibitem{MohajerinEsfahaniKuhn2018} Peyman Mohajerin Esfahani and Daniel Kuhn. \emph{Data-Driven Distributionally Robust Optimization Using the Wasserstein Metric: Performance Guarantees and Tractable Reformulations}. Mathematical Programming, 171(1--2):115--166, 2018. doi:10.1007/s10107-017-1172-1. arXiv:1505.05116.
\bibitem{WangGaoXie2025} Jie Wang, Rui Gao, and Yao Xie. \emph{Sinkhorn Distributionally Robust Optimization}. Operations Research, 2025. doi:10.1287/opre.2023.0294. arXiv:2109.11926.
\bibitem{Girsanov1960} I. V. Girsanov. \emph{On Transforming a Certain Class of Stochastic Processes by Absolutely Continuous Substitution of Measures}. Theory of Probability and Its Applications, 5(3):285--301, 1960. doi:10.1137/1105027.
\bibitem{CameronMartin1945} R. H. Cameron and W. T. Martin. \emph{Transformations of Wiener Integrals under a General Class of Linear Transformations}. Transactions of the American Mathematical Society, 58(2):184--219, 1945.
\bibitem{Yeh1978} J. Yeh. \emph{Transformation of Conditional Wiener Integrals under Translation and the Cameron--Martin Translation Theorem}. Tohoku Mathematical Journal, 30(4):505--515, 1978.
\bibitem{Kakutani1948} Shizuo Kakutani. \emph{On Equivalence of Infinite Product Measures}. Annals of Mathematics, 49(1):214--224, 1948.
\bibitem{ShiDeBortoliCampbellDoucet2023} Yuyang Shi, Valentin De Bortoli, Andrew Campbell, and Arnaud Doucet. \emph{Diffusion Schr\"odinger Bridge Matching}. Advances in Neural Information Processing Systems, 2023. arXiv:2303.16852.
\bibitem{LiuLipmanNickelKarrerTheodorouChen2024} Guan-Horng Liu, Yaron Lipman, Maximilian Nickel, Brian Karrer, Evangelos A. Theodorou, and Ricky T. Q. Chen. \emph{Generalized Schr\"odinger Bridge Matching}. International Conference on Learning Representations, 2024. arXiv:2310.02233.
\bibitem{DomingoEnrichDrozdzalKarrerChen2025} Carles Domingo-Enrich, Michal Drozdzal, Brian Karrer, and Ricky T. Q. Chen. \emph{Adjoint Matching: Fine-Tuning Flow and Diffusion Generative Models with Memoryless Stochastic Optimal Control}. arXiv preprint arXiv:2409.08861, 2025.

\bibitem{Birkhoff1957} Garrett Birkhoff. \emph{Extensions of Jentzsch's Theorem}. Transactions of the American Mathematical Society, 85(1):219--227, 1957.
\bibitem{Sinkhorn1964} Richard Sinkhorn. \emph{A Relationship Between Arbitrary Positive Matrices and Doubly Stochastic Matrices}. Annals of Mathematical Statistics, 35(2):876--879, 1964.
\bibitem{Sinkhorn1967Prescribed} Richard Sinkhorn. \emph{Diagonal Equivalence to Matrices with Prescribed Row and Column Sums}. American Mathematical Monthly, 74(4):402--405, 1967.
\bibitem{Bauer1965} Friedrich L. Bauer. \emph{An Elementary Proof of the Hopf Inequality for Positive Operators}. Numerische Mathematik, 7:331--337, 1965.
\bibitem{Bushell1973} P. J. Bushell. \emph{Hilbert's Metric and Positive Contraction Mappings in a Banach Space}. Archive for Rational Mechanics and Analysis, 52:330--338, 1973.

\end{thebibliography}

\end{document}
